\chapter*{Abstract}
\addcontentsline{toc}{chapter}{Abstract}

Educational institutions that serve large numbers of students manage a
fragmented workflow: syllabi must be structured manually, printed study
material must be converted to machine-readable text, question sets are
authored by hand, and examinations are assembled, scheduled, attempted, and
evaluated across separate systems. This project designs and implements
CatLium EduTech, a multi-tenant Software-as-a-Service platform that automates
this complete lifecycle from raw syllabus and study material to a delivered,
evaluated examination.

The platform is implemented as a modular monolith: a single NestJS
(TypeScript) API that hosts the business domains, an independent Python
stack that provides an OCR service and asynchronous RabbitMQ-based workers
for AI-driven and material-processing workloads, and a Next.js web
application used by institute administrators, teachers, and students. Every
data access is scoped to an institute through a membership-based tenancy
model, with cookie-based session authentication, CSRF protection, and role
based authorization. Content flows through a deterministic pipeline: syllabus
documents are proposed as chapter--topic hierarchies by an AI assistant
under teacher review; PDF and image materials are converted to page-accurate
text by an OCR service; a Material Intelligence engine cleans, structures,
and segments the extracted text and maps it to syllabus entities; and a
question bank feeds paper patterns, question papers, and assessments that
students attempt with server-side deadline enforcement and deterministic
evaluation of objective question types. Study content (notes, summaries,
flashcards, and concepts) is generated by AI workers through an internal
OpenAI-compatible gateway, with schema validation on both worker and API
boundaries.

The system is validated by 188 API unit tests, 118 Python tests across the
OCR and worker services, 15 web tests, and end-to-end harnesses that
exercise the live stack through its public gateway. The report documents the
analysis, design, implementation, testing, and results of the completed
system, and clearly separates implemented scope from deferred future work.